% !TEX encoding = UTF-8
% !TEX TS-program = pdflatex
% !TEX root = ../tesi.tex

%**************************************************************
\chapter{L'azienda}
\label{cap:processi-metodologie}

\section{Profilo}
WebPD s.a.s. é una software house nata nel 2011, con sede a Padova ma con clienti lungo tutto il nord Italia, specializzata nello sviluppo di software gestionali, applicazioni web e mobile (sia Android che iOS) e in servizi di consulenza informatica.\\
\begin{figure}[!h] 
	\centering 
	\includegraphics[width=0.4\columnwidth]{azienda/logo_webpd} 
	\caption{Logo WebPD}
\end{figure}
\\
Il primo prodotto sviluppato da WebPD é il gestionale Martina, un software per piattaforma Windows, sviluppato secondo un'architettura client-server modulare, che consente la gestione del ciclo attivo (vendite) e passivo (acquisti) di un'impresa commerciale. Grazie alla sua struttura modulare, é possibile aggiungere a Martina diverse funzionalità, come ad esempio gestione di magazzino, gestione della contabilità e gestione dello scadenziario.\\\\
Successivamente l'azienda aggiunge al portfolio di servizi offerti anche la realizzazione, grazie alla piattaforma OpenChart, di siti e-commerce. È in questo contesto che nascono TuttoNauticaWeb.com e MiglioNautico.com .\\\\
Nel 2015 il gruppo di agenzie viaggio Primarete Travel s.r.l. avvia una collaborazione con WebPD s.a.s.: nasce così WebPD s.r.l., di cui Primarete detiene il 49\% delle quote. Il frutto di questa collaborazione é l'ammodernamento, continuo ed ancora in corso, dei portali Viaggiregalo.it e  CrociereRegalo.it, di proprietà di Promoter Travel s.r.l., una controllata di Primarete Travel s.r.l. .\\\\

\section{Dominio tecnologico}
L'azienda offre un ampio ventaglio di prodotti, che coinvolgono diverse piattaforme. È opportuno esaminare le tecnologie adoperate raggruppandole nelle seguenti categorie: \textit{DBMS}, \textit{Desktop}, \textit{Mobile} e \textit{Web}.

\subsection{DBMS}
Per lo sviluppo di prodotti software, WebPD si appoggia a due DBMS: \textit{Microsoft SQL Server 2012} e \textit{MySQL Server}. Entrambi questi software sono RDBMS, ma la scelta di utilizzare l'uno o l'altro dipende dalle caratteristiche del progetto, nello specifico: \begin{itemize}
	\item \textbf{Microsoft SQL Server} viene utilizzato nei progetti che prevedono un'ampia mole di dati da gestire e interrogazioni (query) complesse (come, ad esempio, CrociereRegalo.it), in quanto in tali contesti si è dimostrato mediamente più performante di \textit{MySQL};
	\item \textbf{MySQL Server} viene utilizzato il più possibile in caso il progetto non contenga una grossa mole di dati e/o interrogazioni complesse, in quanto non prevede costi di licenza (mentre \textit{SQL Server} si, e anche abbastanza elevati \footcite{site:sql-server-pricing}) ed é multipiattaforma (attualmente disponibile per Linux, Windows e MacOS).
\end{itemize}

\subsection{Desktop}

\subsection{Mobile}

\subsection{Web}
